% sec13_5.tex — Section 13.5: A Signal Problem for a Vibrating String of Finite Length
\section{A Signal Problem for a Vibrating String of Finite Length}
\label{sec:13.5}

In partial differential equations the Laplace inversion integrals that are needed are often not as simple as in Sec.~13.4.  We illustrate this for a vibrating string of length $L$ initially at rest in the horizontal equilibrium position subject to the following time-dependent boundary condition at one end $x = L$:
\begin{align}
\text{PDE:}\quad &
\begin{array}{|c|}
\hline
\displaystyle\pdd{u}{t} = c^2\pdd{u}{x} \\[4pt]
\hline
\end{array}
\label{eq:13.5.1} \\[8pt]
\text{BC:}\quad &
\begin{array}{|c|}
\hline
u(0,t) = 0 \\
u(L,t) = b(t) \\
\hline
\end{array}
\label{eq:13.5.2} \\[8pt]
\text{IC:}\quad &
\begin{array}{|c|}
\hline
u(x,0) = 0 \\[2pt]
\displaystyle\pd{u}{t}(x,0) = 0. \\
\hline
\end{array}
\label{eq:13.5.3}
\end{align}

The zero initial conditions facilitate the use of the Laplace transform in $t$ of $u(x,t)$:
\begin{equation}
\overline{U}(x,s) = \int_0^{\infty} e^{-st}\,u(x,t)\dt.
\label{eq:13.5.4}
\end{equation}
By transforming \eqref{eq:13.5.1}, $\overline{U}(x,s)$ satisfies the ordinary differential equation
\begin{equation}
s^2\overline{U} = c^2\frac{\partial^2\overline{U}}{\partial x^2},
\label{eq:13.5.5}
\end{equation}
subject to the boundary conditions
\begin{align}
\overline{U}(0,s) &= 0 \label{eq:13.5.6}\\
\overline{U}(L,s) &= B(s), \label{eq:13.5.7}
\end{align}
where $B(s)$ is the Laplace transform of $b(t)$.  We can easily determine $\overline{U}(x,s)$:
\begin{equation}
\overline{U}(x,s) = B(s)\,\frac{\sinh(s/c)x}{\sinh(s/c)L}.
\label{eq:13.5.8}
\end{equation}
The convolution theorem implies that
\begin{equation}
u(x,t) = \int_0^t b(t_0)\,f(t-t_0)\,dt_0,
\label{eq:13.5.9}
\end{equation}
where $f(t)$ is the inverse Laplace transform of
\begin{equation}
F(s) = \frac{\sinh(s/c)x}{\sinh(s/c)L}.
\label{eq:13.5.10}
\end{equation}

To obtain this inverse Laplace transform is not straightforward.  One method to obtain the inverse transform of \eqref{eq:13.5.10} is to attempt to use our elementary tables, in which primarily exponentials appear.  We note that
\[
\frac{\sinh(s/c)x}{\sinh(s/c)L}
= \frac{e^{(s/c)x} - e^{-(s/c)x}}{e^{(s/c)L} - e^{-(s/c)L}}
= \frac{e^{(s/c)x} - e^{-(s/c)x}}{e^{(s/c)L}\bigl(1 - e^{-(2s/c)L}\bigr)}.
\]
However, due to the denominator, this cannot be analyzed in a simple way.  Instead, we can introduce an infinite series of exponentials based on the geometric series $[1/(1-x) = 1 + x + x^2 + \cdots$, if $|x| < 1]$:
\begin{align}
F(s) &= \frac{e^{(s/c)x} - e^{-(s/c)x}}{e^{(s/c)L}\bigl(1 - e^{-(2s/c)L}\bigr)} \notag\\
     &= e^{-(s/c)L}\bigl(e^{(s/c)x} - e^{-(s/c)x}\bigr)
        \bigl(1 + e^{-(2L/c)s} + e^{-(4L/c)s} + \cdots\bigr) \notag\\
     &= \sum_{n=0}^{\infty}\left\{
        \exp\!\left[-s\left(\frac{2nL - x + L}{c}\right)\right]
        - \exp\!\left[-s\left(\frac{2nL + x + L}{c}\right)\right]
        \right\}.
\label{eq:13.5.11}
\end{align}
These are all decaying exponentials since $x < L$, and hence each can be inverted using formula (13.2.2n) in Table~13.2.1.  The Laplace transform is a linear combination of
\[
\exp\!\left[-s\left(\frac{2nL \pm x + L}{c}\right)\right]
\qquad (n \ge 0).
\]
The inverse Laplace transform of $F(s)$ is thus a linear combination of Dirac delta functions, $\delta[t - (2nL \pm x + L)/c]$:
\[
f(t) = \sum_{n=0}^{\infty}\left[
\delta\!\left(t - \frac{2nL - x + L}{c}\right)
- \delta\!\left(t - \frac{2nL + x + L}{c}\right)
\right].
\]

Since $f(t-t_0)$ in \eqref{eq:13.5.9} is the influence function for the boundary condition, these Dirac delta functions represent signals whose travel times are $(2nL \pm x + L)/c$ [elapsed from the signal time~$t_0$].  These can be interpreted as direct signals and their reflections off the boundaries $x_0 = 0$ and $x_0 = L$.  Since the nonhomogeneous boundary condition is at $x_0 = L$, we imagine these signals are initiated there (at $t = t_0$).  The signal can travel to $x$ in different ways, as illustrated in Fig.~13.5.1.  The direct signal must travel a distance $L - x$ at velocity~$c$, yielding the retarded time $(L-x)/c$ (corresponding to $n = 0$).  A signal can also arrive at $x$ by additionally making an integral number of complete circuits, an added travel distance of $2Ln$.  The other terms correspond to waves first reflecting off the wall $x_0 = 0$ before impinging on~$x$.  In this case the total travel distance is $L + x + 2nL$ ($n \ge 0$).  Further details of the solution are left as exercises.

\begin{figure}[H]
\centering
\includegraphics[width=0.8\textwidth]{figures/ch13/fig_13_5_1.png}
\caption{Space-time signal paths.}
\label{fig:13.5.1}
\end{figure}

Using Laplace transforms (and inverting in the way described in this section) yields a representation of the solution as an infinite sequence of reflecting waves.  Similar results can be obtained by the method of characteristics (see Chapter~12) or (in some cases) by using the method of images (sequences of infinite space Green's functions).

Alternatively, in subsequent sections, we will describe the use of contour integrals in the complex plane to invert Laplace transforms.  This technique will yield a significantly different representation of the same solution.

%% ── Exercises 13.5 ──────────────────────────────────────────────
\begin{exercises}{13.5}

\item Consider
\begin{align*}
\pdd{u}{t} &= c^2\pdd{u}{x}, \qquad -\infty < x < \infty \\
u(x,0) &= f(x) \\
\pd{u}{x}(x,0) &= g(x).
\end{align*}
Solve using Laplace transforms:
\begin{enumerate}
\item[(a)] if $g(x) = 0$
\item[(b)] if $f(x) = 0$
\end{enumerate}

\item
\begin{enumerate}
\item[(a)] Using the results of this section, invert \eqref{eq:13.5.8} based on the convolution theorem.  Solve for $u(x,t)$.
\item[(b)] Without using the convolution theorem, from \eqref{eq:13.5.8} and \eqref{eq:13.5.11}, determine $u(x,t)$.
\end{enumerate}

\item[\starred 13.5.3.] Solve for $u(x,t)$ using Laplace transforms:
\[
\pdd{u}{t} = c^2\pdd{u}{x}
\]
\[
u(0,t) = 0, \qquad u(x,0) = 0
\]
\[
\pd{u}{x}(L,t) = b(t), \qquad \pd{u}{t}(x,0) = 0.
\]

\item Reconsider Exercise~13.5.3 if instead
\begin{enumerate}
\item[(a)] $\pd{u}{x}(0,t) = 0$ and $u(L,t) = b(t)$
\item[(b)] $\pd{u}{x}(0,t) = 0$ and $\pd{u}{x}(L,t) = b(t)$
\end{enumerate}

\item Solve for $u(x,t)$ using Laplace transforms
\[
\pd{u}{t} = k\pdd{u}{x}, \qquad -\infty < x < \infty
\]
\[
u(x,0) = f(x).
\]
(\emph{Hint}: See Table~13.2.1 of Laplace transforms.)

\item Solve for $u(x,t)$ using Laplace transforms:
\[
\pd{u}{t} = k\pdd{u}{x}
\]
subject to $u(x,0) = f(x)$, $u(0,t) = 0$, and $u(L,t) = 0$.  By what other method(s) can this representation of the solution be obtained?

\end{exercises}
